\section*{Chapitre \ref{ch:iv1:counting-and-symmetry} --- Dénombrement et symétrie}
\markboth{\protect\omTocMark{}{Réponses~: Dénombrement et symétrie}}{\protect\omTocMark{}{Réponses~: Dénombrement et symétrie}}

\begin{solution}{iq:iv1:counting-and-symmetry:1}
Il y a $7! = 5040$ placements, et $3600$ placent les deux à l'écart.

On fixe la chaise d'une personne pour supprimer les rotations, puis on place les sept autres dans l'ordre. Pour le second décompte, on retranche les placements où les deux sont ensemble. On les colle en un bloc~: il reste alors sept objets autour de la table, ce qui donne $6!$ arrangements, et le bloc admet 2 ordres internes. Donc $5040 - 2 \cdot 720 = 3600$.

\iqfollow{Quelle est cette fraction~?}{$5/7$. Vu de l'un des deux, l'autre s'assoit sur l'une des sept chaises au hasard, et deux de ces chaises sont voisines.}
\end{solution}

\begin{solution}{iq:iv1:counting-and-symmetry:2}
Il y a $34\,650$ arrangements, dont $7350$ ne contiennent jamais deux S consécutifs.

Il y a 11 lettres, avec quatre I, quatre S et deux P, si bien que le coefficient multinomial\index{coefficient multinomial} donne $11!/(4!\,4!\,2!) = 34\,650$.

Pour le second décompte, on place les S en dernier. Les sept autres lettres peuvent être disposées de $7!/(4!\,2!) = 105$ façons. Elles laissent 8 intervalles, en comptant les deux extrémités, et les quatre S se répartissent dans quatre intervalles différents, en $\binom{8}{4} = 70$ façons. Le produit vaut $7350$, soit $7/33$ du total.

\iqwrong{Retrancher du total les «~arrangements avec un bloc SS~». Les blocs se chevauchent de bien des façons, et la méthode des intervalles les écarte tous.}
\end{solution}

\begin{solution}{iq:iv1:counting-and-symmetry:3}
Il y a 252 chemins, et 132 évitent le point.

Un chemin est une suite de cinq pas vers la droite et de cinq pas vers le haut~: il y en a donc $\binom{10}{5} = 252$. Ceux qui passent par $(2, 2)$ se scindent en deux parties indépendantes, $\binom{4}{2} = 6$ façons de l'atteindre et $\binom{6}{3} = 20$ façons de poursuivre. Donc 120 passent par ce point, et $252 - 120 = 132$ l'évitent.

\iqfollow{Comment l'implémenteriez-vous pour de nombreux points bloqués~?}{Remplir un tableau par programmation dynamique~: chaque case contient la somme des cases à sa gauche et en dessous, et zéro si elle est bloquée.}
\end{solution}

\begin{solution}{iq:iv1:counting-and-symmetry:4}
Il y a 286 façons si des boîtes peuvent rester vides, et 84 façons sans aucune boîte vide.

On utilise la méthode des étoiles et des barres\index{méthode des étoiles et des barres}. On aligne les 10 pièces et on place 3 cloisons parmi elles. Cela revient à choisir 3 places sur 13, donc $\binom{13}{3} = 286$. Si aucune boîte ne doit être vide, on place d'abord une pièce dans chacune. Il reste alors 6 pièces à répartir librement~: $\binom{9}{3} = 84$.

\iqwrong{$4^{10}$, qui est la réponse pour des pièces étiquetées. Les pièces étant identiques, le résultat ne dépend que du nombre de pièces tombant dans chaque boîte.}
\end{solution}

\begin{solution}{iq:iv1:counting-and-symmetry:5}
Il y a 315 comités. Le raccourci donne 600, car il compte un comité une fois pour chaque femme qui en fait partie.

On compte le complémentaire\index{dénombrement par complémentaire}~: parmi les $\binom{11}{4} = 330$ comités, $\binom{6}{4} = 15$ ne sont composés que d'hommes, et $330 - 15 = 315$.

La méthode fausse désigne une femme comme «~la~» femme, puis choisit les trois autres, ce qui donne $5 \binom{10}{3} = 600$. Un comité comptant $w$ femmes est produit $w$ fois, une fois pour chacune d'elles. Sommer le nombre de femmes sur tous les comités donne exactement 600.

\iqwhy{Chaque fois qu'un décompte semble trop grand, il faut se demander combien de fois chaque objet est produit. Cette question permet de repérer la plupart des erreurs de dénombrement.}
\end{solution}

\begin{solution}{iq:iv1:counting-and-symmetry:6}
Il y a 252 résultats différents, et la probabilité d'un full house est $300/7776 \approx 0.0386$.

Un résultat sans tenir compte de l'ordre est un multiensemble de 5 valeurs prises parmi 6, donc la méthode des étoiles et des barres\index{méthode des étoiles et des barres} donne $\binom{10}{5} = 252$.

Les résultats ne sont pas équiprobables, on ne peut donc pas diviser par 252. Un full, disons 3, 3, 3, 5, 5, provient de $\binom{5}{2} = 10$ lancers ordonnés, tandis que 1, 2, 3, 4, 5 en donne 120 et 6, 6, 6, 6, 6 n'en donne qu'un seul. Il faut donc compter les lancers ordonnés~: 6 choix pour le brelan, 5 pour la paire, 10 ordres, soit 300 sur $6^5 = 7776$.

\iqwrong{Diviser par 252 pour obtenir $30/252 \approx 0.119$. Traiter les résultats non ordonnés comme équiprobables est l'erreur la plus fréquente sur le thème de ce chapitre.}
\end{solution}

\begin{solution}{iq:iv1:counting-and-symmetry:7}
Les deux valent $3/8$.

On utilise la symétrie\index{echangeabilite@échangeabilité}. Chaque ordre des huit boules est équiprobable, donc chaque position a la même chance de contenir une boule donnée. Chacune des 8 boules a autant de chances d'être la dernière, et 3 d'entre elles sont rouges. Il en va de même pour la deuxième position, et pour toute autre.

\iqwrong{Conditionner sur tous les tirages précédents. Cela marche, mais cela prend une page, et la symétrie rend sans objet l'ordre de sortie des boules. La première boule que l'on voit modifie bien la probabilité conditionnelle de la deuxième, mais avant de voir quoi que ce soit, chaque position a la même chance.}
\end{solution}

\begin{solution}{iq:iv1:counting-and-symmetry:8}
Il y a 8 cubes avec trois faces peintes, 96 avec deux, 384 avec une et 512 sans aucune. Un cube unitaire a en moyenne $0.6$ face peinte.

Les coins donnent les 8. Chacune des 12 arêtes compte $10 - 2 = 8$ cubes intérieurs, ce qui donne 96. Chacune des 6 faces compte $8^2 = 64$ carrés intérieurs, ce qui donne 384. L'intérieur est un bloc de $8^3 = 512$. Ces valeurs se somment à 1\,000, ce qui sert de contrôle.

Pour la moyenne, on compte les faces peintes depuis l'extérieur~: le grand cube a 6 faces de 100 carrés unitaires chacune, soit 600 faces. Rapportées aux 1\,000 cubes, cela fait $0.6$.

\iqfollow{Et pour un côté de $n$~?}{$6n^2/n^3 = 6/n$ faces en moyenne.}
\end{solution}

\begin{solution}{iq:iv1:counting-and-symmetry:9}
Avoir exactement 3 numéros gagnants a une probabilité $\approx 0.0177$, soit environ 1 chance sur 57, et en avoir au moins 3 une probabilité $\approx 0.0186$, soit environ 1 chance sur 54.

On compte les tirages selon leur répartition entre vos 6 numéros et les 43 autres~: la loi hypergéométrique\index{loi hypergéométrique}. Il y a $\binom{49}{6} = 13\,983\,816$ tirages. Les tirages ayant exactement 3 numéros en commun sont au nombre de $\binom{6}{3}\binom{43}{3} = 20 \cdot 12\,341 = 246\,820$. En sommant les cas de 3, 4, 5 et 6 numéros en commun, on obtient $0.0186$.

\iqwrong{Utiliser une loi binomiale avec $p = 6/49$. Les tirages se font sans remise, donc les correspondances ne sont pas indépendantes. La loi binomiale donne une valeur voisine de $0.025$, trop élevée.}
\end{solution}

\enlargethispage{-5\baselineskip}%
\begin{solution}{iq:iv1:counting-and-symmetry:10}
Il y en a 266.

On applique l'inclusion-exclusion\index{inclusion-exclusion}~: on compte les multiples de chaque nombre premier et on les retranche, puis on rajoute ceux retranchés deux fois, puis on retire ceux rajoutés trop souvent~:
\[
1000 - (500 + 333 + 200) + (166 + 100 + 66) - 33 = 266.
\]
Les termes intermédiaires sont les multiples de 6, 10 et 15, et le dernier terme compte les multiples de 30.

Une vérification rapide se fait par la densité. Une part $\tfrac12 \cdot \tfrac23 \cdot \tfrac45 = \tfrac{4}{15}$ des entiers n'est divisible par aucun de ces trois nombres premiers, et $1000 \cdot 4/15 \approx 266.7$. Le même motif se répète tous les 30 entiers, avec 8 survivants dans chaque bloc.

\iqwrong{Multiplier $500 \cdot \tfrac23 \cdot \tfrac45$ en arrondissant à chaque étape. Le résultat est proche par hasard, mais un produit arrondi n'est pas un décompte~: l'inclusion-exclusion donne le nombre exact.}
\iqfollow{Comment l'implémenteriez-vous pour de nombreux nombres premiers~?}{Crible jusqu'à $n$, ou parcourir les $2^k$ sous-ensembles de $k$ nombres premiers, en affectant à chacun le signe de son sous-ensemble. Le parcours des sous-ensembles est le meilleur choix lorsque $n$ est très grand et $k$ petit.}
\end{solution}

\begin{solution}{iq:iv1:counting-and-symmetry:11}
La probabilité vaut $150/243 = 50/81 \approx 0.617$.

On applique l'inclusion-exclusion\index{inclusion-exclusion} sur les boîtes vides. Il y a $3^5 = 243$ affectations. Celles qui laissent vide une boîte donnée sont au nombre de $2^5$, et il y a trois boîtes. Celles qui laissent deux boîtes données vides ne sont qu'une seule, et il y a trois paires. Le nombre d'affectations sans boîte vide est donc
\[
3^5 - 3 \cdot 2^5 + 3 \cdot 1 = 150.
\]
Pour vérifier, on compte par formes. Les effectifs des boîtes sont 3, 1, 1 ou 2, 2, 1. Il y a $3 \cdot \binom{5}{3} \cdot 2 = 60$ cas de la première forme et $3 \cdot \binom{5}{2}\binom{3}{2} = 90$ de la seconde, et $60 + 90 = 150$.

\iqwrong{Commencer par placer une boule dans chaque boîte et laisser les autres aller n'importe où~: $5 \cdot 4 \cdot 3 \cdot 3^2 = 540$, ce qui dépasse le total. N'importe quelle boule pouvait être la «~première~» de sa boîte, si bien que chaque affectation est comptée plusieurs fois.}
\iqfollow{Que calcule-t-on d'autre~?}{C'est la même chose que la probabilité qu'un collectionneur de 3 coupons ait réuni tous les coupons en 5 tirages au plus.}
\end{solution}

\begin{solution}{iq:iv1:counting-and-symmetry:12}
Aucune bonne lettre a pour probabilité $44/120 = 11/30 \approx 0.367$, exactement deux bonnes lettres ont pour probabilité $20/120 = 1/6$, et la première tend vers $1/e \approx 0.368$.

Une permutation sans point fixe est un dérangement\index{dérangement}. Par inclusion-exclusion\index{inclusion-exclusion},
\[
D_n = n! \sum_{k=0}^{n} \frac{(-1)^k}{k!}, \quad\text{donc}\quad D_5 = 44, \quad \frac{D_n}{n!} \to e^{-1}.
\]
La convergence est très rapide. Dès cinq lettres, on est à moins de $0.002$ de la limite.

Pour exactement deux bonnes lettres, on choisit lesquelles, puis on dérange les trois autres~: $\binom{5}{2} D_3 = 10 \cdot 2 = 20$.

\iqwrong{Penser que la probabilité qu'aucune lettre ne soit bien placée doit fortement croître ou décroître avec le nombre de lettres. Il y a davantage de lettres à placer, mais chacune a moins de chances d'être bien placée, et les deux effets se compensent~: le nombre espéré de bonnes places vaut toujours un, et le nombre de bonnes places tend vers une loi de Poisson de moyenne 1.}
\iqfollow{Est-il possible que exactement quatre lettres soient bien placées~?}{Non. Si quatre lettres sont bien placées, la cinquième enveloppe est la seule qui reste, et la lettre qui lui correspond doit l'être aussi.}
\end{solution}

\enlargethispage{-3\baselineskip}%
\begin{solution}{iq:iv1:counting-and-symmetry:13}
La probabilité vaut $(6 - 4)/(6 + 4) = 1/5$.

C'est le théorème du scrutin\index{théorème du scrutin}, et la preuve utilise le principe de réflexion\index{principe de réflexion}. Parmi les $\binom{10}{4} = 210$ ordres, on compte les mauvais, ceux où le décompte est à égalité à un certain moment. Un ordre qui commence par un vote pour B est mauvais à coup sûr, puisque A doit rattraper son retard, ce qui conduit à une égalité. Il y a $\binom{9}{3} = 84$ de ces ordres.

Prenons maintenant un ordre qui commence par A et présente une égalité. On échange les votes de A et de B jusqu'à la première égalité. On obtient ainsi un ordre qui commence par B, et l'échange peut être annulé. Les mauvais ordres qui commencent par A sont donc eux aussi au nombre de 84. Les bons ordres sont au nombre de $210 - 168 = 42$, et $42/210 = 1/5$.

\iqwrong{Raisonner comme si «~A a plus de votes, donc il est en général en tête~» et donner un grand nombre. Une avance dès le tout premier vote est une exigence forte, et elle échoue dès que B prend la tête au début.}
\iqfollow{Et avec 60 votes contre 40~?}{La probabilité suit la même formule~: $20/100 = 1/5$. Seule compte la marge rapportée au nombre de votes, ce qui montre que la taille du décompte n'a pas d'importance.}
\end{solution}

\begin{solution}{iq:iv1:counting-and-symmetry:14}
Avec cinq billets de chaque sorte, la probabilité vaut $1/6$. Avec six billets de 5~\$ et quatre billets de 10~\$, elle vaut $3/7$.

On note un billet de 5~\$ comme un pas montant et un billet de 10~\$ comme un pas descendant. La caisse dispose de la monnaie nécessaire exactement lorsque la marche ne passe jamais sous 0. On applique le principe de réflexion\index{principe de réflexion} aux mauvais ordres~: ceux qui touchent $-1$. On réfléchit la partie de la marche jusqu'à la première fois où elle touche $-1$. Elle part alors de $-2$ et finit au même point, si bien que les mauvais ordres se mettent en correspondance avec tous les chemins de $-2$ au point d'arrivée. Pour cinq et cinq, un tel chemin comporte six pas montants et quatre pas descendants, donc il y en a $\binom{10}{4} = 210$, face au total de $\binom{10}{5} = 252$. Il y a donc 42 bons ordres, soit le nombre de Catalan\index{nombres de Catalan} $C_5$, et $42/252 = 1/6$.

Le résultat général pour $a$ billets de 5~\$ et $b$ billets de 10~\$ est $(a + 1 - b)/(a + 1)$. Pour six et quatre, on obtient $3/7$~: $90$ des 210 ordres.

\iqwrong{Supposer qu'avec des billets en nombre égal la probabilité avoisine un demi. Disposer de la monnaie à la toute fin est automatique, et c'est le chemin avant la fin qui échoue.}
\iqwhy{Le principe de réflexion est l'outil qui se cache derrière le prix d'une option à barrière dans un modèle binomial~: un chemin de prix qui franchit un niveau est mis en correspondance avec un chemin réfléchi qui finit ailleurs.}
\end{solution}

\begin{solution}{iq:iv1:counting-and-symmetry:15}
La probabilité vaut exactement $1/2$.

On compare après les dix premiers lancers de chacun. Soit $X$ le nombre de piles du joueur à 11 lancers sur les 10 premiers et $Y$ celui de son adversaire. Par symétrie\index{argument de symétrie}, $P(X > Y) = P(X < Y) = a$, disons, et $P(X = Y) = 1 - 2a$. Ce joueur gagne si $X > Y$, ou si $X = Y$ et que sa dernière pièce tombe sur pile. Si $X < Y$, même un pile final ne donne au plus qu'une égalité, et une égalité n'est pas strictement plus. Sa chance vaut donc $a + \tfrac12(1 - 2a) = \tfrac12$.

Une version plus élégante~: le joueur à 11 lancers a soit plus de piles que son adversaire, soit plus de faces, et jamais les deux, car ses 11 lancers dépassent les 10 de l'adversaire. Pour une pièce équilibrée, pile et face jouent le même rôle, donc chacun a une probabilité $1/2$.

\iqwrong{Penser que la pièce supplémentaire donne plus de la moitié. Elle donne bien le départage, mais les égalités étaient défavorables au départ, si bien que ce gain ramène tout juste à un demi.}
\iqfollow{Et avec une pièce biaisée~?}{La symétrie tombe. Si pile est plus probable que face, le joueur à 11 lancers gagne plus d'une fois sur deux, car l'événement «~plus de piles~» est favorisé.}
\end{solution}

\begin{solution}{iq:iv1:counting-and-symmetry:16}
Elles sont équiprobables~: chacune a une probabilité exactement égale à $1/52$.

On construit le jeu mélangé en deux étapes. On retire la carte qui nous intéresse, et l'on mélange les 51 autres. On insère alors cette carte à une place au hasard, parmi les 52 intervalles possibles\index{argument d'insertion}.

Pour le deux de trèfle~: elle se trouve «~juste après le premier as~» pour exactement un des 52 intervalles, quelle que soit la disposition des 51 autres cartes. Sa probabilité vaut donc $1/52$.

Pour l'as de pique~: on reprend la même donne. Les 51 autres cartes contiennent trois as, et le premier d'entre eux se trouve quelque part. Si l'on insère l'as de pique juste après lui, il devient la carte suivant le premier as. Si on l'insère plus tôt, il devient lui-même le premier as, et il ne suit pas un as. Si on l'insère plus tard, il ne vient pas juste après. Un seul intervalle convient donc, et là encore la probabilité vaut $1/52$.

\iqwrong{Dire que l'as de pique est moins probable, parce qu'il a pu déjà sortir comme premier as, ou plus probable, parce que les as se regroupent. Ces deux intuitions sont justes sur un point, mais le décompte par insertion tranche.}
\iqwhy{Construire un objet aléatoire une carte à la fois à partir d'un objet plus simple évite souvent un conditionnement lourd, et les recruteurs cherchent ce genre de démarche.}
\end{solution}

\begin{solution}{iq:iv1:counting-and-symmetry:17}
La probabilité vaut $\frac{39 \cdot 26 \cdot 13}{51 \cdot 50 \cdot 49} \approx 0.105$.

On place les as un par un, en voyant le jeu comme 52 emplacements répartis entre quatre mains\index{conditionnement séquentiel}. Le premier as peut aller n'importe où. Le deuxième doit éviter les 12 autres emplacements de sa main, il dispose donc de 39 emplacements convenables sur les 51 restants. Le troisième dispose de 26 emplacements convenables sur les 50 restants, et le quatrième de 13 sur les 49 restants.

L'autre voie consiste à compter les donnes entières. Le nombre de donnes où chaque main reçoit un as vaut $4!\binom{48}{12}\binom{36}{12}\binom{24}{12}$, à comparer aux $\binom{52}{13}\binom{39}{13}\binom{26}{13}$ donnes possibles au total. On retrouve la même fraction, après de nombreuses simplifications.

\iqwrong{La réponse indépendante, $4!/4^4 \approx 0.094$. Elle traite chaque as comme choisissant une main de façon indépendante, et elle ignore le fait qu'une main qui contient un as a moins de place pour le suivant.}
\iqfollow{Quelle méthode va le plus loin~?}{Placer les cartes spéciales une à une fonctionne tant qu'elles sont peu nombreuses. Compter les donnes est préférable lorsqu'il faut contrôler plusieurs cartes à la fois.}
\end{solution}

\begin{solution}{iq:iv1:counting-and-symmetry:18}
Les deux sommes comptent chacune 27 issues, et la somme 14 sur quatre dés en compte 146.

En listant les triplets non ordonnés pour 10, on obtient $(1,3,6)$, $(1,4,5)$, $(2,2,6)$, $(2,3,5)$, $(2,4,4)$, $(3,3,4)$, avec $6, 6, 3, 6, 3, 3$ ordres, soit 27. Elles sont égales par une symétrie\index{argument de symétrie} du dé~: remplacer chaque face $x$ par $7 - x$ est une bijection des issues qui envoie une somme $s$ sur $21 - s$, et $21 - 10 = 11$.

Pour quatre dés, on utilise la méthode des étoiles et des barres\index{méthode des étoiles et des barres} avec l'inclusion-exclusion\index{inclusion-exclusion}. En posant $d_i = 1 + y_i$, il faut $y_1 + y_2 + y_3 + y_4 = 10$ avec $y_i \le 5$. Sans plafond, cela donne $\binom{13}{3} = 286$ solutions. Si un $y_i \ge 6$, on lui retire 6 et l'on compte les solutions d'une somme égale à 4~: $\binom{7}{3} = 35$, pour chacun des 4 dés. Deux dés ne peuvent pas dépasser en même temps, car $12 > 10$. Donc $286 - 140 = 146$.

\iqwhy{La méthode des étoiles et des barres avec plafond revient au même calcul qu'une convolution discrète, et la symétrie miroir divise par deux le travail chaque fois qu'on calcule à la main une table de dés.}
\end{solution}

\begin{solution}{iq:iv1:counting-and-symmetry:19}
Il y a $5775$ répartitions, et la probabilité que les deux personnes se retrouvent ensemble vaut $3/11$.

Pour le nombre de répartitions, on part de $12!$ ordres, on découpe en trois blocs de quatre, et l'on divise par les ordres à l'intérieur de chaque bloc et par les ordres des blocs~: $12!/(4!^3\,3!) = 5775$. Autrement, on choisit les 3 compagnons d'une personne fixée, puis les 3 compagnons de la première personne restante~: $\binom{11}{3}\binom{7}{3} = 165 \cdot 35 = 5775$.

La seconde partie ne demande aucun calcul. Fixons la première personne. Son groupe compte 3 autres places, et par symétrie\index{argument de symétrie} la seconde personne a la même chance de se trouver à l'une des 11 autres places. La probabilité vaut donc $3/11$. Le décompte le confirme~: $\binom{10}{2}\binom{7}{3} = 1575$ des 5775 répartitions les réunissent.

\iqwrong{Donner $1/3$ parce qu'il y a trois groupes. Ce serait juste si chaque personne choisissait un groupe de façon indépendante, mais le groupe de la première personne n'a plus que trois places libres, alors que les autres en ont chacun quatre.}
\iqfollow{Avec des groupes étiquetés~?}{Le nombre de répartitions devient $12!/4!^3 = 34\,650$, soit six fois plus. La probabilité reste la même.}
\end{solution}

\begin{solution}{iq:iv1:counting-and-symmetry:20}
Les deux joueurs se rencontrent avec une probabilité $1/32$. Le deuxième meilleur joueur atteint la finale avec une probabilité $32/63$.

Pour la première partie, on compte les matchs plutôt que les tours. À chaque match, un seul joueur est éliminé, donc un tournoi de 64 joueurs comporte 63 matchs. Il y a $\binom{64}{2} = 2016$ paires de joueurs. Par symétrie\index{argument de symétrie}, chaque paire a la même chance d'être la paire d'un match donné, donc chaque paire se rencontre avec une probabilité $63/2016 = 1/32$, et, plus généralement, $2/n$.

Pour la seconde partie, le deuxième meilleur ne peut perdre que contre le meilleur. Les deux se trouvent dans des moitiés opposées du tableau exactement lorsque le deuxième meilleur atteint la finale. Étant donnée la place du meilleur joueur, il reste 63 autres places, dont 32 se trouvent dans l'autre moitié. La probabilité vaut donc $32/63$.

\iqwrong{Parcourir les tours, avec un facteur pour chaque tour que les deux doivent franchir. Cela se fait, mais c'est beaucoup plus long. La linéarité appliquée aux matchs donne la réponse en une ligne.}

\iqwhy{C'est pourquoi le finaliste battu dans un tournoi à élimination directe est une mauvaise estimation du deuxième meilleur~: il a raison un peu plus d'une fois sur deux.}
\end{solution}

\begin{solution}{iq:iv1:counting-and-symmetry:21}
Six points donnent 31 régions, et non 32, et dix points en donnent 256.

La formule est $1 + \binom{n}{2} + \binom{n}{4}$\index{régions d'un disque}. On part du disque, qui forme une seule région. On trace les cordes une à une. Une nouvelle corde ajoute une région pour elle-même, et une de plus pour chaque corde qu'elle croise, car chaque croisement coupe en deux une région déjà existante. Le nombre de régions vaut donc $1 + (\text{cordes}) + (\text{croisements})$. Il y a $\binom{n}{2}$ cordes, et chaque ensemble de 4 points donne exactement un croisement, celui de ses deux diagonales. Il y a donc $\binom{n}{4}$ croisements. Pour $n = 6$~: $1 + 15 + 15 = 31$. Pour $n = 10$~: $1 + 45 + 210 = 256$.

La formule donne $1, 2, 4, 8, 16$ pour $n = 1, \dots, 5$ seulement parce que $\sum_k \binom{n-1}{k}$ coïncide avec elle pour les petites valeurs de $n$. Un polynôme de degré 4 ne peut pas suivre $2^{n-1}$.

\iqwrong{Trente-deux, d'après le motif. Les recruteurs posent cette question pour vérifier si un candidat va extrapoler à partir de quelques cas sans raison.}
\iqfollow{Comment vérifier 31 sans formule~?}{Par la formule d'Euler. Avec $n + \binom{n}{4}$ sommets et $n + \binom{n}{2} + 2\binom{n}{4}$ arêtes, le nombre de faces intérieures vaut $E - V + 1$, ce qui donne le même résultat.}
\end{solution}

\begin{solution}{iq:iv1:counting-and-symmetry:22}
Il y a 292 façons.

On utilise la programmation dynamique\index{comptage de monnaie} sur les pièces, une dénomination à la fois. Soit $w(a)$ le nombre de façons de former le montant $a$ avec les pièces utilisées jusqu'ici, en partant de $w(0) = 1$. Pour chaque pièce $c$, à tour de rôle, on fait parcourir $a$ à partir de $c$ vers le haut et on pose $w(a) \leftarrow w(a) + w(a - c)$. Parcourir vers le haut permet de réutiliser une pièce, et traiter les pièces une par une ne compte chaque combinaison qu'une seule fois, quel que soit l'ordre. C'est la même chose que lire un coefficient de la fonction génératrice $\prod_c 1/(1 - x^c)$.

À la main, on raisonne sur le nombre de grosses pièces. Sans pièce de 50 et sans pièce de 25, par exemple, choisir $d$ pièces de 10 laisse $21 - 2d$ choix de pièces de 5, et les pièces de 1 complètent le reste, ce qui donne 121 façons pour ce cas. En sommant tous les cas, on obtient 292. Une boucle exhaustive sur les pièces de 50, 25, 10 et 5, les pièces de 1 complétant le reste, donne le même nombre.

\iqwrong{Parcourir d'abord les montants, puis les pièces à l'intérieur. Cela compte $1 + 5$ et $5 + 1$ comme deux cas différents, c'est-à-dire qu'on compte des ordres au lieu de combinaisons.}
\iqfollow{Et le nombre de façons en tenant compte de l'ordre~?}{Il suffit d'échanger les deux boucles. Il devient énorme, car la plupart des suites utilisent beaucoup de pièces de 1.}
\end{solution}

\begin{solution}{iq:iv1:counting-and-symmetry:23}
Il y a 136 solutions.

Sans plafond, les étoiles et les barres\index{méthode des étoiles et des barres} donnent $\binom{32}{2} = 496$. Il faut maintenant retrancher les solutions dans lesquelles une variable dépasse 15. Si $x \ge 16$, on pose $x' = x - 16$. On a alors $x' + y + z = 14$, qui admet $\binom{16}{2} = 120$ solutions. Il y a trois cas de ce type, et deux variables ne peuvent pas être toutes deux au moins égales à 16, puisque leur somme dépasserait 30. L'inclusion--exclusion\index{inclusion--exclusion} s'arrête donc là~: $496 - 360 = 136$.

Une vérification géométrique~: les solutions forment une tranche triangulaire d'un cube, et avec un plafond égal à la moitié du total, la tranche est le triangle central formé par les trois milieux. Son aire vaut environ un quart de celle du triangle non contraint, et $136/496 \approx 0.27$, un peu plus à cause des points du bord.

\iqwhy{Les plafonds apparaissent dès qu'une capacité limite un partage. C'est la manière standard de les traiter, et elle s'arrête après un ou deux termes lorsque le plafond est grand.}
\end{solution}

\begin{solution}{iq:iv1:counting-and-symmetry:24}
La première probabilité vaut $1/70$, et la seconde $1/34\,650$.

On ignore toutes les autres cartes. Seul compte l'ordre relatif des 8 as et rois, et toutes les façons de les placer sont également probables\index{argument de symétrie}. Parmi les $\binom{8}{4} = 70$ motifs de A et de K, un seul a tous les A en premier. Les 44 autres cartes ne jouent aucun rôle.

Pour douze cartes dans l'ordre as, rois, dames, on regarde les motifs à trois lettres~: il y en a $12!/(4!)^3 = 34\,650$, et un seul est le bon.

\iqwrong{Multiplier les probabilités conditionnelles pour chaque as, par exemple $\tfrac{4}{52}$ pour le premier, puis essayer de suivre le jeu. Retirer les cartes inutiles est précisément l'idée.}

\iqwhy{Ce type de réduction, où seul compte l'ordre relatif de quelques événements, s'applique à toute question sur l'événement qui survient en premier parmi des instants échangeables, comme l'ordre de remplissage dans une suite d'arrivées aléatoires.}
\end{solution}

\begin{solution}{iq:iv1:counting-and-symmetry:25}
La probabilité vaut $\binom{20}{10}/4^{10} = 184\,756/1\,048\,576 \approx 0.176$.

La somme directe vaut $\sum_k \binom{10}{k}^2 / 4^{10}$. La méthode élégante voit l'événement «~nombres égaux de piles~» comme un énoncé sur 20 lancers. On modifie le sens de mes lancers en échangeant pile et face, ce qui les laisse équilibrés. Alors «~vos piles égalent mes piles~» devient «~vos piles plus mes nouvelles piles font 10~», car mes nouvelles piles valent $10$ moins mes anciennes piles. C'est l'événement «~exactement 10 piles sur 20 lancers équilibrés~», dont la probabilité vaut $\binom{20}{10}/2^{20}$.

Le même raisonnement démontre l'identité de Vandermonde\index{identité de Vandermonde} $\sum_k \binom{n}{k}^2 = \binom{2n}{n}$, et pour $n$ grand la probabilité vaut environ $1/\sqrt{\pi n}$, soit ici $0.178$.

\iqwrong{Un sur 11, puisqu'il y a 11 valeurs possibles du nombre de piles. Ces valeurs ne sont pas équiprobables, et les valeurs centrales portent l'essentiel de la masse.}
\iqwhy{La probabilité d'une égalité exacte après $n$ tours ne décroît qu'avec la racine carrée de $n$. Les égalités dans un grand échantillon sont bien plus fréquentes que l'intuition ne le laisse penser.}
\end{solution}

\begin{solution}{iq:iv1:counting-and-symmetry:26}
L'inclusion a pour probabilité $(3/4)^{10} \approx 0.056$, la disjonction a la même probabilité, et l'égalité a pour probabilité $(1/2)^{10} = 1/1024$.

Chaque élément agit indépendamment\index{comptage élément par élément}, et il y a quatre cas pour chacun~: dans $A$ seulement, dans $B$ seulement, dans les deux, ou dans aucun. Pour $A \subseteq B$, seul le cas «~dans $A$ seulement~» est interdit, donc chaque élément contribue $3/4$. Il y a donc $3^{10}$ couples favorables parmi les $4^{10}$.

Pour la disjonction, seul le cas «~dans les deux~» est interdit, et là encore chaque élément contribue $3/4$. Passer au complémentaire de $B$ transforme l'un en l'autre, ce qui explique l'égalité des deux probabilités. Pour l'égalité, un élément doit être soit dans les deux, soit dans aucun, ce qui a pour probabilité $1/2$.

\iqwrong{Conditionner sur la taille de $A$, puis sommer sur cette taille. Cela donne bien la même réponse, par le théorème du binôme, mais la vision élément par élément évite toute somme.}
\iqfollow{Et pour $A \subseteq B \subseteq C$~?}{Chaque élément a 8 cas, dont 4 sont permis~: $(1/2)^{10}$.}
\end{solution}

\enlargethispage{-3\baselineskip}%
\begin{solution}{iq:iv1:counting-and-symmetry:27}
Dans une rangée il y a 816 choix, et sur un cercle il y en a 800.

Dans une rangée, on utilise un argument sur les écarts\index{bijection}. On retire une chaise entre chaque paire voisine de chaises choisies. Choisir 3 chaises parmi 20 sans deux voisines revient alors à choisir 3 chaises parmi $20 - 2 = 18$, sans condition. Donc $\binom{18}{3} = 816$.

Sur un cercle, on distingue selon que la chaise 1 est choisie. Si elle l'est, ses deux voisines sont exclues, et il faut 2 chaises non voisines parmi les 17 autres, disposées en rangée~: $\binom{16}{2} = 120$. Sinon, il faut 3 chaises non voisines parmi 19 disposées en rangée~: $\binom{17}{3} = 680$. Le total vaut 800. La formule générale est $\frac{n}{n-k}\binom{n-k}{k}$.

La rangée et le cercle diffèrent par les sélections qui contiennent les deux chaises d'extrémité, lesquelles ne sont voisines que sur le cercle. En fixant ces deux chaises, la troisième doit venir des chaises 3 à 18, soit 16 possibilités, et $816 - 16 = 800$ exactement.

\iqwhy{La différence entre une rangée et un anneau est un piège récurrent. Pour vérifier une réponse circulaire, il est commode de compter ce que la fermeture de la rangée en anneau supprime.}
\end{solution}

\begin{solution}{iq:iv1:counting-and-symmetry:28}
Le compte vaut $1\,263\,360$ sur $10! = 3\,628\,800$, soit environ $0.348$.

On utilise l'inclusion--exclusion\index{inclusion--exclusion} sur l'ensemble des couples assis ensemble. On fixe $k$ couples et on colle chacun en un bloc. Il reste alors $10 - k$ objets, qui peuvent être disposés de $(10 - k)!$ façons, et chaque bloc a 2 ordres internes, donc le terme vaut $\binom{5}{k} 2^k (10 - k)!$. D'où
\[
\sum_{k=0}^{5} (-1)^k \binom{5}{k} 2^k (10 - k)! = 3\,628\,800 - 3\,628\,800 + 1\,612\,800 - 403\,200 + 57\,600 - 3840.
\]
Les deux premiers termes s'annulent exactement, ce qui est une petite surprise, et la somme vaut $1\,263\,360$. La même formule donne 8 et 240 pour deux et trois couples, ce que l'on vérifie en les énumérant.

Une approximation sert de contrôle utile. Le nombre espéré de couples ensemble vaut $5 \cdot 2/10 = 1$, et si ce nombre suivait une loi de Poisson\index{approximation de Poisson} de moyenne 1, la réponse serait $e^{-1} \approx 0.368$. La vraie réponse est un peu plus basse, car les événements ne sont pas indépendants, et elle tend vers $1/e$ pour un grand nombre de couples.

\iqwrong{Multiplier $(1 - 1/5)^5 \approx 0.328$ comme si les événements étaient indépendants. C'est une estimation honnête, mais elle n'est pas exacte, alors qu'un compte exact était demandé.}
\iqfollow{Comment une table ronde modifie-t-elle le résultat~?}{Des blocs collés disposés sur un cercle donnent $(9 - k)!$ arrangements avec les mêmes signes, et le compte diffère. Ne pas oublier le facteur des rotations.}
\end{solution}

\begin{solution}{iq:iv1:counting-and-symmetry:29}
Il y a 208 telles parties.

On utilise le filtre des racines de l'unité\index{filtre des racines de l'unité}. Soit $\omega = e^{2\pi i/5}$, et considérons $f(x) = \prod_{k=1}^{10}(1 + x^k)$. Le coefficient de $x^s$ compte les parties de somme $s$. La moyenne de $f$ sur les cinq racines cinquièmes de l'unité ne conserve que les termes pour lesquels $s$ est divisible par 5~:
\[
N = \frac15 \sum_{j=0}^{4} f(\omega^j).
\]
Pour $j = 0$, on a $f(1) = 2^{10} = 1024$. Pour $j \ne 0$, les puissances $\omega^{jk}$ parcourent deux fois toutes les racines cinquièmes de l'unité lorsque $k$ va de 1 à 10, donc $f(\omega^j) = \bigl[\prod_{r=0}^{4}(1 + \omega^r)\bigr]^2$. Comme $\prod_r (x - \omega^r) = x^5 - 1$, en posant $x = -1$ on obtient $\prod_r (1 + \omega^r) = 2$. Chacun des quatre autres termes vaut donc $2^2 = 4$, et
\[
N = \frac{1024 + 4 \cdot 4}{5} = 208.
\]
Les autres résidus se partagent le reste de façon égale, soit 204 pour chaque résidu non nul. La répartition est donc presque uniforme, et le léger excédent pour le résidu 0 est exactement ce que détecte le filtre.

\iqwrong{Répondre $1024/5$, qui n'est pas un entier, ou supposer une uniformité parfaite. Chaque bloc $\{1, \dots, 5\}$ contient un multiple de 5, donc les résidus ne sont pas parfaitement équilibrés.}
\iqfollow{Et pour les parties de $\{1, \dots, 5\}$~?}{Les puissances ne parcourent qu'une fois les racines, donc $(32 + 4 \cdot 2)/5 = 8$. L'énumération le confirme.}
\end{solution}

\begin{solution}{iq:iv1:counting-and-symmetry:30}
Il y a 4 colliers à rotation près, et 3 si le retournement est aussi permis.

On utilise le lemme de Burnside\index{lemme de Burnside}~: le nombre de colliers distincts est la moyenne, sur les symétries, du nombre de motifs de perles laissés inchangés par chaque symétrie. Il y a $\binom{6}{3} = 20$ motifs.

Sous la rotation, l'identité laisse inchangés les 20 motifs. Une rotation d'un ou de cinq crans déplace toutes les perles dans un seul cycle, ce qui suppose qu'elles aient toutes la même couleur~: 0 motif. Une rotation de deux ou de quatre crans a deux cycles de longueur 3, donc l'un doit être rouge et l'autre vert~: 2 motifs chacune. Une rotation de trois crans a trois cycles de longueur 2, et le nombre de perles rouges devrait être pair~: 0. La moyenne vaut $(20 + 2 + 2)/6 = 4$. Ce sont RRRGGG, RRGRGG, RRGGRG et RGRGRG.

En ajoutant les retournements, les trois réflexions passant par deux perles opposées laissent chacune 4 motifs, car les deux perles situées sur l'axe doivent être de couleurs différentes et l'une des deux autres paires est rouge. Les trois réflexions passant par des espaces entre perles n'en laissent aucun, car les paires rendent le nombre de perles rouges pair. La moyenne vaut $(24 + 12)/12 = 3$. Les deux images miroir RRGRGG et RRGGRG se réduisent à une seule.

\iqwrong{Diviser $20$ par les 6 rotations, ce qui donne $10/3$. Un motif qui possède une symétrie a une orbite plus petite, donc la simple division est fausse, et le lemme sert précisément à la corriger.}
\end{solution}

\begin{solution}{iq:iv1:counting-and-symmetry:31}
Il y a 10 cubes à deux couleurs, et 57 à trois couleurs.

On utilise le lemme de Burnside\index{lemme de Burnside} sur les 24 rotations du cube. Une rotation laisse un coloriage inchangé exactement lorsque chaque cycle de faces qu'elle crée est d'une seule couleur, de sorte qu'avec $k$ couleurs, une rotation à $c$ cycles laisse inchangés $k^c$ coloriages. Les rotations se répartissent selon leur nombre de cycles~:
\begin{itemize}
\item l'identité a 6 cycles (une rotation)~;
\item un demi-tour autour de l'axe passant par deux faces opposées a 4 cycles (trois rotations)~;
\item un quart de tour autour de cet axe a 3 cycles (six rotations)~;
\item un demi-tour autour de l'axe passant par deux arêtes opposées a 3 cycles (six rotations)~;
\item une rotation autour de l'axe passant par deux sommets opposés a 2 cycles (huit rotations).
\end{itemize}
Le total fait $1 + 3 + 6 + 6 + 8 = 24$, et le nombre de cubes est
\[
\frac{k^6 + 3k^4 + 12k^3 + 8k^2}{24},
\]
ce qui donne 10 pour $k = 2$ et 57 pour $k = 3$. Pour $k = 1$, on obtient 1, un contrôle utile sur le compte des éléments du groupe.

Pour $k = 2$, on peut énumérer selon le nombre de faces de couleur A, de 0 à 6~: $1, 1, 2, 2, 2, 1, 1$. Ces valeurs font 10, ce qui confirme la formule.

\iqwrong{$2^6/24$, qui n'est même pas un entier. Les coloriages ayant une symétrie ont des orbites plus petites.}
\iqfollow{Et ceux qui utilisent chacune des six couleurs une seule fois~?}{Aucune rotation autre que l'identité n'en laisse un inchangé, donc $6!/24 = 30$.}
\end{solution}

\begin{solution}{iq:iv1:counting-and-symmetry:32}
Un angle droit a pour probabilité $60/220 = 3/11$. Le fait de contenir strictement le centre a pour probabilité $40/220 = 2/11$.

Il y a $\binom{12}{3} = 220$ triangles. Par l'angle inscrit dans un demi-cercle\index{angle inscrit dans un demi-cercle}, un triangle dont les sommets sont sur un cercle a un angle droit exactement lorsqu'un de ses côtés est un diamètre. Il y a 6 diamètres et 10 choix pour le troisième sommet, ce qui donne 60.

On parcourt le polygone et on écrit les trois arcs entre sommets consécutifs sous forme d'entiers dont la somme vaut 12. Un triangle contient strictement le centre exactement lorsque chaque arc est inférieur à 6. Un arc égal à 6 correspond au cas du triangle rectangle, et un arc supérieur à 6 laisse le centre à l'extérieur du triangle (triangle obtus).

On compte les triangles obtus. On fixe le sommet où se termine le plus grand arc, de sorte que les deux autres sommets se trouvent dans les 5 pas suivants~: il y a 12 choix pour ce sommet et $\binom{5}{2} = 10$ choix pour les deux autres. Il y a donc 120 triangles obtus, et $220 - 60 - 120 = 40$ triangles aigus.

\iqwrong{Utiliser $1/4$, qui est la réponse pour trois points uniformes sur un cercle. Un polygone fini a des triangles rectangles qui prennent une part de la probabilité, donc la part des triangles aigus est plus faible. Quand le nombre de sommets augmente, la réponse tend bien vers $1/4$.}
\iqfollow{Trois points uniformes~: pourquoi $1/4$~?}{Les arcs forment un point uniforme sur un simplexe, et la région où tous les arcs sont inférieurs à la moitié est le triangle central, qui vaut un quart du simplexe.}
\end{solution}

\begin{solution}{iq:iv1:counting-and-symmetry:33}
On utilise le principe des tiroirs\index{principe des tiroirs}, en comparant le nombre de sous-ensembles au nombre de sommes possibles. Pour 7 nombres, le nombre de sous-ensembles est trop petit pour forcer une coïncidence, et les puissances 1, 2, 4, 8, 16, 32, 64 ont leurs 127 sommes de sous-ensembles toutes distinctes.

Tout ensemble de 10 nombres compris entre 1 et 100 possède $2^{10} - 1 = 1023$ sous-ensembles non vides. Chacun a une somme comprise entre 1 et au plus $91 + 92 + \dots + 100 = 955$. Il y a plus de sous-ensembles que de sommes possibles, donc deux sous-ensembles différents ont la même somme.

Ces deux sous-ensembles peuvent se chevaucher. On retire leurs éléments communs aux deux~: les sommes restent égales entre elles, et les ensembles restent différents. Aucun ne devient vide, car les nombres sont strictement positifs, et si l'un d'eux était vide, l'autre devrait avoir une somme nulle. On obtient donc deux groupes non vides et disjoints de même somme.

Pour 7 nombres, il n'y a que 127 sous-ensembles, et les sommes peuvent atteindre 679, si bien que l'argument des tiroirs ne donne rien. Les puissances de 2 en sont un contre-exemple, car les chiffres binaires sont uniques. L'énoncé lui-même est donc faux à 7, et pas seulement sa preuve.

\iqwrong{Chercher à construire les deux groupes. L'argument prouve qu'ils existent sans montrer lesquels ils sont, et c'est précisément l'intérêt de la question.}
\iqwhy{Les arguments des tiroirs donnent des bornes très rapides, et la même comparaison de comptages fournit une borne inférieure sur le nombre de bits qu'exige une table de hachage.}
\end{solution}

\begin{solution}{iq:iv1:counting-and-symmetry:34}
La probabilité vaut $232/1024 = 29/128 \approx 0.227$.

On utilise le principe de réflexion\index{principe de réflexion}. Un chemin qui atteint $+4$ et finit en dessous de 4 peut être réfléchi après son premier contact avec 4, ce qui donne un chemin finissant au-dessus de 4. C'est une bijection. Donc
\[
P(\max \ge 4) = P(S_{10} \ge 4) + P(S_{10} > 4).
\]
Les chemins qui finissent à 4 ou au-dessus doivent comporter au moins 7 pas montants~: $\binom{10}{7} + \binom{10}{8} + \binom{10}{9} + \binom{10}{10} = 120 + 45 + 10 + 1 = 176$. Les chemins finissant strictement au-dessus de 4 en comportent au moins 8~: 56. Le total vaut $176 + 56 = 232$ sur $2^{10}$.

\iqwrong{Doubler $P(S_{10} \ge 4)$, pour obtenir $352/1024$. La formule en temps continu donne bien un facteur deux, mais ici les chemins qui finissent exactement en 4 doivent être comptés une seule fois, et non deux fois. Sur seulement 10 pas, la différence est considérable.}
\iqfollow{Que vous dit la limite continue~?}{Le mouvement brownien sur une durée $T$ atteint un niveau $a$ avec probabilité $2P(W_T > a)$, où les contacts exacts n'ont pas de masse. L'écart diminue avec le nombre de pas.}
\iqwhy{C'est la réflexion utilisée pour les options à barrière~: la probabilité de toucher un niveau pendant une période vaut environ le double de la probabilité de finir au-delà. Près du niveau, sur une grille grossière, la correction discrète compte.}
\end{solution}

\enlargethispage{-3\baselineskip}%
\begin{solution}{iq:iv1:counting-and-symmetry:35}
Il y a $55\,252$ tickets porte-bonheur.

L'astuce tient à une bijection\index{bijection}. On remplace chacun des trois derniers chiffres $d$ par $9 - d$. Un ticket est porte-bonheur lorsque $a + b + c = d + e + f$. Après le remplacement, $d' + e' + f' = 27 - (d + e + f)$, donc la condition devient $a + b + c + d' + e' + f' = 27$. On compte donc les chaînes de six chiffres dont la somme vaut 27, ce qui se ramène à un seul coefficient~: celui de $x^{27}$ dans $(1 + x + \dots + x^9)^6$.

On calcule ce nombre par la méthode des étoiles et des barres\index{méthode des étoiles et des barres} avec l'inclusion--exclusion\index{inclusion--exclusion}. Sans plafond, le nombre de solutions vaut $\binom{32}{5} = 201\,376$. On retranche 6 fois le nombre de cas où un chiffre vaut au moins 10, soit $\binom{22}{5} = 26\,334$, et on ajoute 15 fois le nombre de cas où deux chiffres sont dans ce cas, soit $\binom{12}{5} = 792$. Trois chiffres supérieurs ou égaux à 10 demanderaient une somme d'au moins 30, donc on s'arrête là~:
\[
201\,376 - 158\,004 + 11\,880 = 55\,252.
\]
On vérifie par la somme directe des carrés~: pour chaque valeur $s$ de 0 à 27, on élève au carré le nombre de chaînes de trois chiffres de somme $s$, puis on additionne ces carrés~; on retrouve le même total. Le terme maximal vaut $75^2$, pour $s = 13$.

\iqwrong{Supposer qu'environ un ticket sur 28 est porte-bonheur, soit 35\,714. Les sommes proches du centre sont beaucoup plus probables, donc une coïncidence est bien plus probable que cela~: la vraie proportion est d'environ 1 sur 18.}
\end{solution}

\begin{solution}{iq:iv1:counting-and-symmetry:36}
Il existe $2^{16} = 65\,536$ tableaux de ce type.

On raisonne en termes de degrés de liberté\index{degrés de liberté}. On remplit le bloc $4 \times 4$ en haut à gauche de n'importe quelle manière, ce qui fait $2^{16}$ choix. Chacune des quatre premières lignes voit alors sa dernière case imposée pour que la somme de la ligne soit paire, et il en va de même pour les quatre premières colonnes. Il ne reste que la case du coin inférieur droit.

Le coin doit fixer à la fois la dernière ligne et la dernière colonne, et ces deux exigences s'accordent. La somme des quatre premières cases de la dernière colonne et la somme des quatre premières cases de la dernière ligne ont chacune la parité de la somme du bloc $4 \times 4$~: chacune s'obtient en additionnant les lignes, ou les colonnes, du bloc. Le coin existe donc et il est unique, quel que soit le bloc choisi.

En général, le compte pour une grille $n \times n$ vaut $2^{(n-1)^2}$. L'énumération donne 2, 16 et 512 pour $n = 2, 3, 4$, comme il se doit.

\iqwrong{Compter naïvement équations et inconnues~: 25 inconnues, moins 10 contraintes, laissent $2^{15}$. Les 10 équations ne sont pas indépendantes, car la somme de toutes les lignes et la somme de toutes les colonnes coïncident, si bien qu'une équation est redondante.}
\iqwhy{Les grilles de parité sont l'idée derrière le stockage RAID et les codes correcteurs d'erreurs simples. L'équation redondante est ce qui permet de retrouver une case défaillante.}
\end{solution}

\enlargethispage{-3\baselineskip}%
\begin{solution}{iq:iv1:counting-and-symmetry:37}
Le compte vaut $2^n - n - 1$, soit 26 pour $n = 5$.

Une permutation ayant exactement une montée est constituée de deux séries décroissantes placées l'une après l'autre\index{nombres eulériens}. Chacune est déterminée par son ensemble d'éléments, car une série décroissante n'admet qu'un seul ordre. On choisit donc un sous-ensemble $S$ pour la première série~: il y a $2^n$ sous-ensembles, et les autres éléments forment la seconde série.

Chaque permutation ayant une montée provient d'un unique sous-ensemble. Mais certains sous-ensembles ne donnent aucune montée~: si le dernier élément de la première série est plus grand que le premier de la seconde, toute la permutation descend et il n'y a pas de montée. Cela arrive exactement lorsque $S$ est l'ensemble des $k$ plus grandes valeurs, pour $k = 0, 1, \dots, n$~: cela fait $n + 1$ sous-ensembles, chacun donnant l'unique permutation entièrement décroissante. Le compte vaut donc $2^n - (n + 1)$.

L'énumération le confirme~: $n = 3$ donne 4, à savoir 132, 213, 231 et 312, et $n = 4$ donne 11, $n = 5$ donne 26, $n = 6$ donne 57.

\iqwrong{Choisir où placer la montée, parmi $n - 1$ positions, puis un sous-ensemble de taille fixée. Cela ne donne le bon total qu'à condition de retrancher, pour chaque taille, les sous-ensembles qui ne donnent aucune montée, et c'est un chemin plus long vers le même compte.}
\iqfollow{Comment appelez-vous ces nombres~?}{Des nombres eulériens. Ils comptent les permutations selon leur nombre de montées, et ils apparaissent lorsqu'on additionne les puissances des entiers.}
\end{solution}

\begin{solution}{iq:iv1:counting-and-symmetry:38}
La paire est $\{1, 2, 2, 3, 3, 4\}$ et $\{1, 3, 4, 5, 6, 8\}$, connue sous le nom de dés de Sicherman\index{dés de Sicherman}, et c'est la seule autre paire.

On utilise des fonctions génératrices\index{fonctions génératrices}. Un dé est un polynôme $D(x) = \sum x^{\text{face}}$, et la somme de deux dés a pour fonction génératrice le produit de leurs fonctions. La cible est
\[
(x + x^2 + \dots + x^6)^2 = x^2 (1 + x)^2 (1 + x + x^2)^2 (1 - x + x^2)^2,
\]
et ces facteurs sont irréductibles sur les entiers. La factorisation polynomiale est unique, donc chaque dé $D$ est un produit de certains de ces facteurs. Deux conditions s'appliquent à chaque dé. D'abord, il doit vérifier $D(0) = 0$, car les faces sont positives, donc il contient un facteur $x$. Ensuite, il doit vérifier $D(1) = 6$, car il possède six faces. En $x = 1$, les facteurs valent $2$, $3$ et $1$. Pour obtenir 6, chaque dé doit prendre exactement un $(1 + x)$ et un $(1 + x + x^2)$.

Seuls les facteurs $(1 - x + x^2)$, dont la valeur en 1 vaut 1, peuvent se déplacer. En les répartissant $1 + 1$, on obtient les dés standard. En les répartissant $2 + 0$, on obtient
\[
x(1 + x)(1 + x + x^2) = x + 2x^2 + 2x^3 + x^4
\]
et $x(1 + x)(1 + x + x^2)(1 - x + x^2)^2 = x + x^3 + x^4 + x^5 + x^6 + x^8$.

Une recherche informatique sur tous les ensembles de faces le confirme~: il n'existe exactement que ces deux solutions.

\iqwrong{Chercher à la main en imposant «~les deux dés doivent totaliser 42 ensemble~». C'est une condition nécessaire vraie, mais elle ne mène pas à la paire, et elle ne prouve rien sur l'unicité.}
\iqfollow{Et avec trois dés~?}{La même factorisation offre davantage de possibilités, et il existe des triplets non standard, également obtenus en déplaçant les facteurs $(1 - x + x^2)$.}
\end{solution}

\begin{solution}{iq:iv1:counting-and-symmetry:39}
Il y a 312 placements, et la probabilité qu'un placement alterné aléatoire convienne vaut $13/120 \approx 0.108$.

Il s'agit du problème des ménages\index{problème des ménages}. Fixer un homme permet de supprimer les rotations~: les autres hommes peuvent être placés de $4! = 24$ façons. Entre les hommes, il y a cinq places pour les femmes, et chaque femme est interdite sur deux places voisines~: celles situées à côté de son partenaire. Il faut donc compter les affectations de 5 femmes à 5 places, avec ces positions interdites, et ce nombre ne dépend pas de la façon dont les hommes sont assis.

Les interdictions forment un seul anneau de 10 cases~: femme, place, femme, place, et ainsi de suite. Choisir $k$ violations revient à choisir $k$ cases de l'anneau sans que deux soient voisines, car une femme ou une place ne peut pas servir deux fois. Il existe $\frac{2n}{2n - k}\binom{2n - k}{k}$ choix de ce type, et les $n - k$ autres sont placées librement, donc, par inclusion-exclusion\index{inclusion-exclusion},
\[
M_n = \sum_{k=0}^{n} (-1)^k \frac{2n}{2n - k}\binom{2n - k}{k} (n - k)!, \qquad M_5 = 13.
\]
Le total vaut $4! \cdot 13 = 312$, sur $4! \cdot 5! = 2880$ placements alternés. Un calcul exhaustif sur les 2\,880 placements par machine donne le même compte, et la formule donne aussi $M_6 = 80$.

\iqwrong{Traiter chaque femme comme ayant indépendamment une chance de $3/5$ d'obtenir une place valide, ce qui donne $(3/5)^5 \approx 0.078$. Les choix se disputent les places, et ils sont loin d'être indépendants.}
\iqfollow{Et pour $n$ grand~?}{Chaque femme a 2 places interdites sur $n$, donc le nombre attendu de violations vaut 2, et la fraction tend vers $e^{-2} \approx 0.135$, ce qui est une bonne vérification ici.}
\end{solution}

\begin{solution}{iq:iv1:counting-and-symmetry:40}
Chaque prisonnier doit commencer par la boîte portant son propre numéro, puis suivre les numéros qu'elle contient. La probabilité de succès vaut $1 - \sum_{k=51}^{100} 1/k \approx 0.312$, à comparer avec $2^{-100}$ pour un choix au hasard.

Chaque boîte contient un numéro qui est l'étiquette d'une autre boîte, donc les boîtes forment une permutation\index{structure en cycles des permutations}. Si un prisonnier ouvre la boîte $i$, trouve le numéro $j$, puis ouvre la boîte $j$ et ainsi de suite, il parcourt le cycle qui contient $i$. Ce cycle se termine à la boîte qui contient $i$, ce que le prisonnier cherche exactement. Le prisonnier réussit donc si et seulement si le cycle est de longueur au plus 50. Tous réussissent si et seulement si la permutation n'a aucun cycle de longueur supérieure à 50.

Une permutation de 100 éléments a au plus un cycle de longueur supérieure à 50. Le nombre de permutations ayant un cycle de longueur exactement $k > 50$ se calcule en choisissant ses éléments, $\binom{100}{k}$, en les disposant en cycle, $(k - 1)!$, puis en disposant le reste, $(100 - k)!$. Cela vaut $100!/k$. Ainsi la probabilité d'un cycle long vaut $\sum_{k=51}^{100} 1/k \approx \ln 2$, et le succès a une probabilité d'environ $1 - \ln 2 \approx 0.307$.

La stratégie fonctionne parce que le succès est fortement corrélé~: tous les prisonniers réussissent ou échouent ensemble, selon une seule propriété de la permutation.

\iqwrong{«~Chacun a $1/2$, donc la probabilité vaut $2^{-100}$ quoi qu'ils fassent, car personne ne peut transmettre d'information.~» La probabilité individuelle vaut bien $1/2$. Ce que le plan modifie, c'est la corrélation entre eux, et c'est précisément l'indépendance qu'on abandonne.}
\iqwhy{C'est le meilleur exemple d'un schéma qui compte pour le risque~: on ne peut pas modifier le comportement marginal d'une position, mais c'est la corrélation qui détermine ce qui arrive à l'ensemble des positions.}
\end{solution}
